(1) 整个系统刚性且隔热，去掉隔板不做边界功；两室气体初温相同，理想气体内能只依赖温度，因此总内能守恒给终温仍为 \(T\)。总摩尔数为 \(n=n_1+n_2\)、总容积 \(V=nRT/p\)，所以终压强仍为 \(p\)。按照 Dalton 分压定律，组分 \(i\) 的分压为 \(p_i=n_iRT/V=x_i p\)。

(2) 对任一理想气体组分，以其可用体积从 \(V_i\) 扩至 \(V\) 的等温可逆参考过程计算熵变：
\[
\Delta S_i=n_iR\ln\frac V{V_i}=-n_iR\ln x_i.
\]
故总混合熵
\[
\Delta S_{\rm mix}=-R(n_1\ln x_1+n_2\ln x_2)>0,
\]
因为 \(0<x_i<1\)。真实移除隔板过程不可逆，但熵是状态函数，允许用虚拟可逆路径求初末差，不能反过来断言真实过程可逆。

(3) 理想混合物组分化学势满足
\[
\mu_i(T,p_i)=\mu_i^\circ(T)+RT\ln(p_i/p^\circ).
\]
因各组分压力由初始 \(p\) 变为 \(p_i=x_ip\)，
\[
\Delta G_{\rm mix}=RT\sum_i n_i\ln x_i=-T\Delta S_{\rm mix}<0.
\]
理想地使用各组分专用半透膜进行反向可逆等温等压分离时，最小外界功为 \(W_{\rm sep,min}=-\Delta G_{\rm mix}=T\Delta S_{\rm mix}\)；真实设备需要更多功。

(4) 若 \(n_1=n_2=n_0\)，则 \(x_1=x_2=1/2\)，所以 \(\Delta S_{\rm mix}=2n_0R\ln2\)，\(W_{\rm sep,min}=2n_0RT\ln2\)。但若隔板两侧装的是同一种、同温同压的理想气体，去除隔板前后是不可区分的宏观平衡态，真正的混合熵为零。错误的 \(2n_0R\ln2\) 来自将两批本来不可区分的同种分子人为标记为两种可用半透膜区分的组分，属于微观态重复计数；不能仅靠对每半区套“自由膨胀”公式消除。
